% Irodov, Part II section 2.3 final problems 104--112, printed pp.74--75
% (PDF pp.83--84). All formulas use classical Boltzmann distributions.
\thermokineticproblem{104}{74}
在地表，氢分子与氮分子的数密度比值为
$\eta_0=n_{\mathrm H_2}(0)/n_{\mathrm N_2}(0)$。
若温度 $T=280\,\mathrm K$、
自由落体加速度 $g$ 在高度范围内均恒定，
求高度 $h=3000\,\mathrm m$ 处的相同比值
$\eta$ 与 $\eta_0$ 的比。
\exsolution
等温玻尔兹曼分布给
\[
n_i(h)=n_i(0)
\exp\left(-\frac{M_i gh}{RT}\right),
\]
其中 $M_i$ 是气体组分摩尔质量。
因此
\[
\eta(h)=\frac{n_{\mathrm H_2}(h)}
{n_{\mathrm N_2}(h)}
=\eta_0\exp\left[
\frac{(M_{\mathrm N_2}-M_{\mathrm H_2})gh}{RT}
\right].
\]
所以
\[
\boxed{\frac{\eta}{\eta_0}=
\exp\left[
\frac{(M_{\mathrm N_2}-M_{\mathrm H_2})gh}{RT}
\right]}.
\]
取 $M_{\mathrm N_2}=0.028\,\mathrm{kg/mol}$、
$M_{\mathrm H_2}=0.002\,\mathrm{kg/mol}$，
\[
\boxed{\eta/\eta_0\simeq1.39}.
\]
\excheck 较轻的氢分子随高度衰减得慢，
因此高处氢/氮数密度比上升。

\thermokineticproblem{105}{74}
一根很高的竖直容器里充有两类分子，
质量分别为 $m_1<m_2$。
底部数密度分别为 $n_1<n_2$。
假设温度 $T$、重力加速度 $g$ 都与高度无关，
求两类分子数密度相等的高度 $h$。
\exsolution
对每种粒子应用玻尔兹曼分布：
\[
n_1(h)=n_1e^{-m_1gh/(k_BT)},\qquad
n_2(h)=n_2e^{-m_2gh/(k_BT)}.
\]
令二者相等，
\[
\ln(n_2/n_1)=
\frac{(m_2-m_1)gh}{k_BT}.
\]
所以
\[
\boxed{h=\frac{k_BT\ln(n_2/n_1)}
{(m_2-m_1)g}}.
\]
\excheck 由于分子质量与底部浓度都满足严格大小关系，
得到 $h>0$；较重组分浓度随高度下降更快，
所以存在唯一相等的高度。

\thermokineticproblem{106}{74}
一个\emph{足够高}的竖直圆柱容器
内装一定质量的二氧化碳，
温度原为 $T$。
假设气体等温分布、
重力场均匀，
若气体升温至 $\eta T$，
底部压强将怎样变化？
\exsolution
取容器底面积为 $S$、
顶部高度可近似取无穷远，且顶部压强趋于零。
竖直静力平衡为
$\mathrm dp/\mathrm dh=-\rho(h)g$。
两边从底部 $h=0$ 积分到无穷：
\[
p(\infty)-p(0)
=-g\int_0^\infty\rho(h)\,\mathrm dh.
\]
因为 $p(\infty)=0$，
总气体质量为
$m_{\rm gas}=S\int_0^\infty\rho(h)\,\mathrm dh$，
故
\[
\boxed{p(0)=\frac{m_{\rm gas}g}{S}}.
\]
加热会让气体重新分布到更高位置，
但总质量和底面积不变，所以
\[
\boxed{p_{\rm bottom,new}=p_{\rm bottom,old}}.
\]
\excheck 高度必须足以使顶面压强可忽略；
若容器不够高、顶部有非零外压，
结论需要重新考虑边界条件。

\thermokineticproblem{107}{74}
一个足够高的竖直圆柱容器
装有温度为 $T$ 的理想气体。
在均匀重力场中，
求每个气体分子的平均重力势能；
这一平均值是否取决于气体是否由不同种类分子组成？
\exsolution
对单一分子质量 $m$ 的组分，
玻尔兹曼分布沿高度为
$n(h)\propto e^{-mgh/(k_BT)}$。
平均高度：
\[
\bar h=\frac{\int_0^\infty h e^{-h/H}\,\mathrm dh}
{\int_0^\infty e^{-h/H}\,\mathrm dh}
=H=\frac{k_BT}{mg}.
\]
因此单分子的平均重力势能
\[
\boxed{\overline U=mg\bar h=k_BT}.
\]
若混合气含多种分子，每种组分 $i$ 的
平均高度虽为
$H_i=k_BT/(m_i g)$，
但其平均势能均满足
$m_i gH_i=k_BT$。
故任意组分数目加权平均仍为
\[
\boxed{\overline U_{\rm mixture}=k_BT}.
\]
\excheck 混合气不同分子在高度上的分布不同，
但在这种理想化的无限高等温重力场中，
各组分平均重力势能相同。

\thermokineticproblem{108}{74}
一根长 $\ell=100\,\mathrm{cm}$、
两端密封的水平管道，
沿自身轴向做恒定加速度 $a$。
管内装有温度为 $T=330\,\mathrm K$ 的氩气，
在管随动参考系中达到平衡。
求使两端氩气数密度相差
$\eta=1.0\%$ 的加速度大小。
\exsolution
在管的加速参考系中，
气体分子受到恒定惯性力，
其沿管道的势能差为
$\Delta U=ma\ell$。
两端数密度的玻尔兹曼比值
\[
\frac{n_{\max}}{n_{\min}}
=e^{ma\ell/(k_BT)}.
\]
题目说差异为 $\eta=1\%$，
可取该比值为 $1+\eta$，
于是
\[
\ln(1+\eta)=\frac{ma\ell}{k_BT}
=\frac{Ma\ell}{RT}.
\]
解得
\[
\boxed{a=\frac{RT}{M\ell}\ln(1+\eta)}.
\]
用氩气 $M=0.040\,\mathrm{kg/mol}$、
$\eta=0.010$：
\[
\boxed{a\simeq6.8\times10^2\,\mathrm{m/s^2}
\simeq70g}.
\]
\excheck 以 $\ln(1.01)\simeq0.01$
只是小百分比近似；精确公式使用自然对数。

\thermokineticproblem{109}{74}
胶体粒子密度为 $\rho$，
溶剂密度为 $\rho_0$。
以恒角速度 $\omega$ 绕竖直轴离心，
同一水平面上两个半径 $r_1,r_2$ 处的
胶体粒子数密度之比为
$n(r_2)/n(r_1)=\eta$。
求胶体粒子每摩尔的质量 $M$。
\exsolution
设单个胶体粒子质量为 $m$、
体积为 $V_p=m/\rho$。
考虑溶剂的浮力类比，
单粒子的有效离心力为
\[
F_{\rm eff}(r)=(m-\rho_0V_p)\omega^2r
=m\left(1-\frac{\rho_0}{\rho}\right)\omega^2r.
\]
有效势能由
$-\mathrm dU/\mathrm dr=F_{\rm eff}$ 得
\[
U(r)=-\frac12m\left(1-\frac{\rho_0}{\rho}\right)
\omega^2r^2+\text{常数}.
\]
用玻尔兹曼分布，
\[
\ln\eta=\frac{m(1-\rho_0/\rho)\omega^2
(r_2^2-r_1^2)}{2k_BT}.
\]
因为 $M=N_Am$、$R=N_Ak_B$，
\[
\boxed{M=\frac{2RT\,\rho\ln\eta}
{(\rho-\rho_0)\omega^2(r_2^2-r_1^2)}}.
\]
\excheck 需要 $(\rho-\rho_0)(r_2^2-r_1^2)\ne0$。
若粒子比溶剂更轻，
其浓度可能向轴心方向增大，
须保留差值符号。

\thermokineticproblem{110}{74--75}
一个两端封闭的水平管道长
$\ell=100\,\mathrm{cm}$，
绕穿过一端的竖直轴以恒角速度 $\omega$ 旋转。
管内装有温度
$T=300\,\mathrm K$ 的二氧化碳，
求使管两端数密度比值为 $\eta=2.0$
的角速度。
\exsolution
在随管旋转的参考系中，
分子的径向离心势为
$U(r)=-m\omega^2r^2/2$。
由玻尔兹曼分布，
\[
n(r)=n(0)\exp\left(\frac{m\omega^2r^2}{2k_BT}\right).
\]
因此
$\eta=n(\ell)/n(0)
=\exp[m\omega^2\ell^2/(2k_BT)]$。
取对数求角速度：
\[
\boxed{\omega=\sqrt{\frac{2k_BT\ln\eta}{m\ell^2}}
=\sqrt{\frac{2RT\ln\eta}{M\ell^2}}}.
\]
用 $M_{\mathrm{CO_2}}\simeq0.044\,\mathrm{kg/mol}$、
$\ell=1.00\,\mathrm m$、$T=300\,\mathrm K$：
\[
\boxed{\omega\simeq2.8\times10^2\,\mathrm{rad/s}}.
\]
\excheck 本版 PDF 的数学抽取可能丢失根号；
$\omega$ 必须由 $\omega^2$ 开平方，
否则不能保证角速度量纲。

\thermokineticproblem{111}{75}
理想气体分子处于某中心场，
单分子势能为 $U(r)=ar^2$，
其中 $a>0$。
温度为 $T$，
场中心数密度为 $n_0$。
求：(a) 位于半径 $[r,r+\mathrm dr]$ 的分子数；
(b) 最概然距场心半径；
(c) 一分子落入该球壳的概率；
(d) 温度降低 $\eta$ 倍且粒子总数固定时，
场心浓度变成原来的多少倍？
\exsolution
玻尔兹曼空间数密度为
$n(r)=n_0e^{-ar^2/(k_BT)}$。
(a) 球壳体积 $4\pi r^2\mathrm dr$，因此
\[
\boxed{\mathrm dN=4\pi n_0r^2
e^{-ar^2/(k_BT)}\,\mathrm dr}.
\]
(b) 径向概率密度正比
$r^2e^{-ar^2/(k_BT)}$。
令其对数导数为零：
$2/r-2ar/(k_BT)=0$，得
\[
\boxed{r_p=\sqrt{\frac{k_BT}{a}}}.
\]
(c) 归一化总数
\[
N=4\pi n_0\int_0^\infty
r^2e^{-ar^2/(k_BT)}\,\mathrm dr
=n_0\pi^{3/2}\left(\frac{k_BT}{a}\right)^{3/2}.
\]
故归一化球壳概率为
\[
\boxed{\frac{\mathrm dN}{N}=
4\pi\left(\frac{a}{\pi k_BT}\right)^{3/2}
r^2e^{-ar^2/(k_BT)}\,\mathrm dr}.
\]
(d) $N$ 固定时 $n_0\propto T^{-3/2}$，
所以 $T\to T/\eta$ 后
\[
\boxed{n_0'/n_0=\eta^{3/2}}.
\]
\excheck 此中心场是三维谐势，
要找\emph{半径}的众数必须包括球壳体积
的 $r^2$ 因子；中心的体积数密度最大
不表示最概然半径为零。

\thermokineticproblem{112}{75}
承接第111题，
分子处于中心势
$U(r)=ar^2$、温度为 $T$。
求：(a) 分子势能位于
$[U,U+\mathrm dU]$ 之间的数目；
(b) 最概然势能 $U_p$，
并与最概然半径 $r_p$ 对应的势能比较。
\exsolution
第111题得到半径球壳分子数
\[
\mathrm dN=4\pi n_0r^2
e^{-ar^2/(k_BT)}\,\mathrm dr.
\]
用变量变换
$U=ar^2$，所以
$r=\sqrt{U/a}$、
$\mathrm dr=\mathrm dU/(2\sqrt{aU})$，
于是
\[
\boxed{\mathrm dN=
\frac{2\pi n_0}{a^{3/2}}
\sqrt U\,e^{-U/(k_BT)}\,\mathrm dU}.
\]
为了找 $U$ 的众数，
对 $\sqrt Ue^{-U/(k_BT)}$ 取对数求导：
\[
\frac1{2U}-\frac1{k_BT}=0
\quad\Longrightarrow\quad
\boxed{U_p=\frac12k_BT}.
\]
而最概然半径
$r_p=\sqrt{k_BT/a}$，
它对应的势能为
\[
\boxed{U(r_p)=ar_p^2=k_BT=2U_p}.
\]
\excheck 半径分布与势能分布各自存在不同的雅可比，
其众数不要求在非线性变换下相互对应。
